\runninghead{Delivered and Waiting}
\properlane{already-and-awaited}{Delivered Already, Waiting Still: The Tenses of a People in Exile}

The texts of this Mass speak of deliverance in two tenses at once. The
Introit's antiphon is all future: \latin{invocábitis \dots\ exáudiam \dots\
redúcam}, ``you shall call \dots\ I will hear \dots\ I will bring back.'' Its
psalm verse is past: \latin{avertísti captivitátem Iacob}, ``Thou hast turned
away the captivity of Jacob.'' The Gradual is past, \latin{Liberásti nos},
``Thou hast delivered us,'' and its verse future, \latin{laudábimur}, ``we will
glory.'' The Epistle ends with converts who wait, \latin{exspectáre Fílium eius
de cælis}, for a Jesus \latin{qui erípuit nos ab ira ventúra}, ``who hath
delivered us from the wrath to come'': the rescue is past and the wrath is
still to come. In the Gospel a seed is sown and becomes a tree, and leaven is
hidden in meal until the whole is leavened. The Communion
promises \latin{accipiétis}, ``you shall receive.'' The question is how a
deliverance can be spoken of as done while those who sing of it are still
waiting for it. St Augustine answers that what God has foretold is as
certain as what has already happened, and St Thomas Aquinas that the cause of
the wrath, sin, has already been taken away. St Jerome asks a related
question, why the exiles must pray for what God has promised, and answers that
God urges them to pray so that they may deserve to receive it.

\subsection{The future spoken of as past}
\label{sec:awaited-augustine}

Augustine meets the question on the Gradual's own verse. ``Thou hast saved us
from our enemies'': ``This too is spoken of the future under the figure of the
past. But this is the reason that it is spoken of as if it were past, that it
is as certain as if it were past.'' He gives the precedent of the Passion,
foretold in the psalm in past tenses, ``They pierced My hands and My feet,''
``They parted My garments among them,'' and not ``They shall pierce'' or
``shall part'': ``All these things are expressed as if they were past,
although they were yet to come: because to God things to come also are as
certain as if they were past.'' Then, on the next verse, the Gradual's ``In God
will we boast all the day long'': ``Observe how he intermingles words
expressive of a future time'' (\work{En.\ in Ps.}~43.6--7). The Gradual sings
exactly the pair of verses in which Augustine finds the past standing for a
certain future and the future set beside it.

On the Introit's verse he makes a second point. ``Thou hast turned away the
captivity of Jacob'' is said to Christ, and the captivity is not to
barbarians: ``Thou hast turned away our captivity, not by setting us free from
the barbarians, with whom we had not met, but by setting us free from bad
works, from our sins, by which Satan held sway over us.'' The psalm goes on,
in past tenses, ``Thou hast forgiven the iniquity of thy people,'' and then
prays, ``Convert us, O God our saviour'' (Ps 84:3, 5). Augustine asks how
the psalmist can pray for what he has just called done, and gives the
psalmist's answer: ``These things I speak of as done, because I see them about
to be done: but because they are not yet done, I pray that they may come,
which I have already seen'' (\work{En.\ in Ps.}~84.2--4).

\subsection{Prayer urged for what is promised}
\label{sec:awaited-jerome}

The Introit's antiphon comes from Jeremias's letter to exiles who were told
to build houses in Babylon and wait seventy years. Jerome, commenting on the
verses the antiphon sings, first says whose knowledge the promise rests on:
\latin{Futurorum igitur scientia soli Deo competit}, ``knowledge of the future
belongs to God alone.'' Then he asks why the exiles must pray at all, since
God has promised: \latin{Certe absque invocatione et oratione captivorum,
poterat implere Dominus quod promiserat; sed hortatur eos ad preces, ut quod
promissum est, mereantur accipere}, ``Surely the Lord could have fulfilled
what he had promised without the calling and prayer of the captives; but he
urges them to prayer, so that they may deserve to receive what was promised.''
He quotes the Lord's \latin{Petite, et accipietis}, ``Ask, and you shall
receive'' (Mt 7:7; \work{In Hieremiam}~V, on 29:10--14, PL~24 col.~893).

Jerome reads the letter anagogically of the Church, and his anagogy has an
order. \latin{Tamdiu in saeculi hujus confusione versamur, quamdiu requiem septenarii
numeri mereamur accipere, et accepta poenitentia, Deus impleat quod promisit,
et reducamur ad locum nostrum Ecclesiam}: ``we remain in the confusion of
this world until we deserve to receive the rest of the number seven; and then,
repentance having been received, God fulfils what he promised, and we are
brought back to our place, the Church'' (col.~893). He continues:
\latin{Idcirco enim Dominus percutere visus est, ut sanaret: dabitque nobis finem laboris nostri
atque patientiae: et invocabimus eum, et revertemur in Ecclesiam, et orabimus,
et exaudiemur}, ``For the Lord seemed to strike in order to heal: and he will
give us the end of our labour and patience: and we shall call upon him, and
return to the Church, and pray, and be heard'' (col.~894). ``The end of our
labour and patience'' is his reading of the clause the antiphon leaves out,
``to give you an end and patience''; ``we shall call upon him \dots\ and be
heard'' is his reading of the antiphon's own words. He then reports, as the
view of others (\latin{Quidam}), that the seventy years are the span of a
life, after which we return to the Lord, \latin{nunc enim in umbra nos omnia
et in imagine possidere}, ``for now we possess all things in shadow and in
image'' (col.~894).

Augustine and Jerome answer related questions from the two parts of the
Introit, its verse and its antiphon. Augustine's psalmist prays because what
he has seen is not yet done; Jerome's exiles pray so that they may deserve to
receive what God, who could have given it without their prayer, has promised.

\subsection{Waiting for the Son from heaven}
\label{sec:awaited-epistle}

The Epistle's last verse has both tenses in one sentence. The Thessalonians
turned from idols \latin{servíre Deo vivo et vero, et exspectáre Fílium eius
de cælis}, ``to serve the living and true God, And to wait for his Son from
heaven \dots\ Jesus, who hath delivered us from the wrath to come.'' St John
Chrysostom, on the verse, sets out what is waited for: ``to wait for
His Son from heaven; Him that was crucified, Him that was buried; to wait for
Him from heaven \dots\ You see all things at the same time; both the
Resurrection, and the Ascension, and the second Coming, the Judgment, the
retribution of the just, the punishment of the wicked.'' He names the
condition of those who wait: ``The terrible things are in hand, but the good
things are in the future, when Christ shall come from heaven. See how much
hope is required'' (\work{Hom.\ in 1~Thess.}~2). His English translator gives
``which delivereth us,'' following the Greek; the Missal's Latin is the
perfect \latin{erípuit}.

St Thomas Aquinas explains that perfect. We still await the resurrection and
deliverance from the punishment to come; but from its cause we are already
freed: \latin{A causa autem poenae, scilicet a peccato, liberamur per
Christum. Unde dicit qui eripuit nos}, ``From the cause of the punishment,
namely from sin, we are freed through Christ. Therefore he says, who
delivered us'' (\work{Super I ad Thessalonicenses} c.~1 l.~1). On the Gradual's
\latin{Salvasti}, the Vulgate's verb, he gives the two times side by side:
God saved us \latin{vel in praesenti: quia etsi corporaliter nobis noceant,
spiritualiter tamen nocere non possunt; vel in futuro}, ``either in the
present, because though they harm us in body they cannot harm us in spirit;
or in the future'' (\work{Super Psalmos}~43, n.~3).

Augustine, Aquinas and Chrysostom put the weight in different places.
Augustine stresses the certainty of what God has foretold; Aquinas, the
removal of the cause now and the escape from the penalty later; Chrysostom,
the hope that the present terrors require.

Paul also remembers \latin{sustinéntiæ spei}, ``the enduring of the hope''
(1~Thess 1:3). Aquinas says that hope \latin{facit patienter sustinere
adversa}, ``makes us bear adversity patiently,'' \latin{ante Deum, non ante
oculos hominum}, ``before God, not before the eyes of men.'' The Collect
describes those who pray as \latin{semper rationabília meditántes}, ``always
meditating on reasonable things.'' For people who, as Paul says, endure in
hope, that \latin{semper} covers the whole time of waiting. The Secret asks
that the offering \latin{gubérnet et prótegat}, ``govern and protect,'' those
who offer it; in this interpretation that keeping, too, is asked for the time
of waiting.

\subsection{A kingdom hidden from the foundation of the world}
\label{sec:awaited-gospel}

Said of the kingdom now, each parable's beginning has happened, the seed sown
and the leaven \latin{abscóndit}, hidden, and its end, the tree and the whole
leavened, is still to come: the same two times as in the chants and the
Epistle. St Peter Chrysologus asks the question the leaven parable's
\latin{donec} raises: \latin{fermentatur totum. Quando? Adveniente
Christo}, ``the whole is leavened. When? When Christ comes'' (1~Cor 13:9--10;
\work{Sermo}~99).

Matthew's citation says what Jesus utters in parables: \latin{eructábo abscóndita
a constitutióne mundi}, ``I will utter things hidden from the foundation of
the world.'' Jerome, on the verse, says that the psalm is spoken \latin{ex
persona Domini}, ``in the person of the Lord,'' and that it narrates the signs
of the exodus. He draws from this a rule for reading: \latin{Ex quo
intelligimus universa illa quae scripta sunt, parabolice sentienda: nec
manifestam tantum sonare litteram, sed et abscondita sacramenta}, ``From this
we understand that all those things that are written are to be taken as
parables: that the letter does not sound only what is plain, but hidden
mysteries as well'' (\work{In Matthaeum}~II, on 13:35, PL~26 col.~93).
Augustine, on the psalm, says that its title calls for a reader who perceives ``some inward sense,''
that ``the mystery of the Kingdom of Heaven was veiled in the Old Testament,
which in the fulness of time should be unveiled in the New,'' and that
parables and propositions are not to be read cursorily (\work{En.\ in
Ps.}~77.1--2). Both, then, read the psalm's account of Israel's past as
carrying more than its letter: for Augustine the mystery veiled in the Old
Testament is unveiled in the New, and for Jerome the letter that tells of the
exodus also sounds hidden mysteries. Jerome's rule is the one this
interpretation applies to the Gradual's psalm and the Introit's prophecy in
answer to the strongest objection, below.

\subsection{The chants as the voice of those who wait}
\label{sec:awaited-chants}

The cry from the depths, sung at the Alleluia and twice at the Offertory, is
the prayer of people who are still in the depths. St Hilary of Poitiers asks
why we must pray at all to a God who is everywhere and hears everything. God
hears, sees and acts, \latin{sed orandus a nobis est, ut secundum precem
nostram adsit, audiat, videat, efficiat}, ``but he is to be prayed to by us,
so that he may be present, hear, see and act according to our prayer.'' And
he puts the answer in one sentence: \latin{Naturae suae est ut audiat; sed
fidei nostrae est ut precemur audiri}, ``It is of his nature that he hears;
but it is of our faith that we pray to be heard'' (\work{Tractatus in
Ps.}~129.2--3). Jerome's question about the exiles, why they must pray for
what God has promised, is related to Hilary's. For Jerome, prayer makes the
exiles worthy to receive what is promised; for Hilary, praying to be heard is
the part faith plays.

The medieval commentators on the Mass who expound these four chants read them
as one sequence. Rupert of Deutz takes the Introit's \latin{cogitationes
pacis} as the promise of the banquet of grace to Christ's brethren according
to the flesh, prefigured in Joseph, who received the brothers who had sold him; in the
Gradual they give thanks for their release \latin{captivitate animarum},
``from the captivity of souls''; in the Offertory they ask the true Joseph,
\latin{Obsecramus ut obliviscaris sceleris fratrum}, ``We beseech thee to
forget the wickedness of thy brethren'' (Gen 50:17); and the Communion is his answer, \latin{Nolite timere \dots\ Deus vertit illud in bonum}, ``Fear not
\dots\ God turned it into good'' (Gen 50:19--21; \work{De divinis
officiis}~XII.23). Sicard of Cremona, without Joseph, reads the same
four chants as the remnant of Israel answering the promise announced in the
Introit, giving thanks in the Gradual and praying in the Offertory, and in the
Communion the Lord \latin{certificat de salute}, ``gives assurance of
salvation'' (\work{Mitrale}~VIII.23). William Durandus takes up Rupert's
Joseph (\work{Rationale}~VI.140). For both Rupert and Sicard the sequence
belongs to the remnant of Israel at the end of time: a promise, thanks for
release, a plea, and the Lord's assurance. They put the Gradual's thanks in
that future although its \latin{Liberásti} is a perfect, as is the
Offertory's \latin{clamávi}, and so, like Augustine, take the Gradual's past
tense of something still to come, though neither argues from the tense as he
does.

\subsection{Believe that you shall receive}
\label{sec:awaited-communion}

The Communion gives a present command and a future gift: \latin{crédite quia
accipiétis}, ``believe that you shall receive.'' Jerome, on the Introit's
verses, quoted a saying of the Lord in the same form, an imperative and a
future, \latin{Petite, et accipietis}; Hilary says that it belongs to faith to
pray to be heard. Neither comments on Mark. The antiphon has all three things
they speak of, asking, believing and receiving, in one sentence, and with
Hilary's sentence in mind its \latin{crédite} can be heard as a call to the faith that
does the asking.

The antiphon names one condition, belief. St Bede, on Mark, adds another for
those whose faith is not yet perfect: they should ask that their sins be
forgiven, \latin{et absque dubio impetrabunt quod petunt, si tamen in se
peccantibus aliis primo ipsi dimittunt}, ``and they will without doubt obtain
what they ask, provided that they first forgive others who sin against them''
(\work{In Marcum}~III, on 11:22--25). Mark's next verse gives the command to
forgive (Mk 11:25).

The Postcommunion is said by people who have received. Its participle is
past, \latin{pasti}, ``having been fed''; its verb is a continuing desire,
\latin{semper \dots\ appetámus}, ``that we may always long for.'' As in the
Epistle's last verse, where Jesus has delivered and the Thessalonians wait,
the meal has been given and the longing for the things by which we truly live
continues.

\subsection{The strongest objection}
\label{sec:awaited-objection}

The plainest reading of these texts gives their past tenses a past. Psalm~43
remembers how God planted his people in the land and how he gave them
victory; its ``Thou hast saved us'' can be read as history, and the psalm
then complains of a defeat. Jeremias's promise was made to particular exiles
and was kept when they returned from Babylon. On that reading the tenses are
not prophecy, and the Mass's past is Israel's past.

Jerome grants that history. He reads the letter first as a letter to
captives, and then anagogically of the Church's return, in an order in which
the confusion of this world lasts until the rest is given. The rule he states
at Mt 13:35 allows both, since by it the letter sounds what is plain as well
as hidden mysteries. Augustine hears the Gradual's psalm as ``the voice of the
Martyrs'' because Paul quoted it to a Church suffering persecution
(\work{En.\ in Ps.}~43.1); for him the psalm's speaker is that Church, and its ``Thou hast saved us'' is a deliverance still
to come, spoken as past because it is certain (43.6).

\begin{fourSenses}
\item[Literal.] Jeremias writes to the exiles in Babylon that God's thoughts
toward them are of peace and that he will bring them back, and tells them to
pray (Jer 29:4--14); Paul recalls that the Thessalonians turned from idols to
wait for God's Son from heaven, Jesus, who hath delivered us from the wrath to
come (1~Thess 1:9--10); Jesus speaks to the crowds in parables, uttering, as
the psalm says,
things hidden from the foundation of the world (Mt 13:34--35).
\item[Allegorical.] The captivity is sin, and Christ has turned it away
(Augustine on Ps 84); the psalm of the exodus is the Lord's own parable, and
all that is written contains hidden mysteries (Jerome, on Mt 13:35); the place
to which the exiles are brought back is the Church (Jerome on Jer 29).
\item[Moral.] Pray for what God has promised: he could give it without our
prayer, and urges us to pray so that we may deserve to receive it (Jerome);
pray, since praying to be heard is faith's part (Hilary), and believe that you
receive (Mk 11:24); forgive, that
your prayer be heard (Bede); endure in hope before God and not before men
(Aquinas on 1~Thess 1:3).
\item[Anagogical.] The Son awaited from heaven, at whose coming the whole is
leavened (Chrysostom on 1~Thess 1:10; Chrysologus); after the confusion of this
world, repentance received, the return to our place, where God, who struck to
heal, gives the end of our labour and patience, and we call and are heard
(Jerome, \work{In Hieremiam}~V, cols.~893--894); the praise of ``all the day,''
which Aquinas takes for eternity (\work{Super Psalmos}~43, n.~4).
\end{fourSenses}
